\section{Năm điều kiện đặt lên một dụng cụ đo}
\label{sec:ch18-nam-dieu-kien}

\index{dụng cụ đo tương lai!năm điều kiện|(}%
Chưa ai chịu trả giá để đo, chứ không phải không ai đo được: đó là chỗ mục trước
dừng. Mục này hỏi cái được mua bằng giá ấy phải có hình dạng nào, và câu trả lời
không phải mới: năm chương liền, mỗi chương đọc một thứ khác nhau, mỗi chương rút
ra một điều kiện rồi giao lại. Mục này xếp cả năm cạnh nhau lần đầu.

\index{dụng cụ đo tương lai!năm điều kiện không bài nào phát biểu}%
Trước khi xếp thì phải nói một điều về nguồn gốc của chúng, vì điều ấy đổi cách
đọc cả bảng. Không điều kiện nào trong năm được phát biểu bởi một bài báo nào của
sách. Mỗi điều kiện là cái còn lại sau khi một chương đọc xong một bài và hỏi bài
ấy chưa nói gì. Nên bảng~\ref{tab:ch18-nam-dieu-kien} là bảng của sách, và nó là
phần trả lời cho nửa sau câu hỏi mà chương này mang tên.

\begin{table}[htbp]
  \centering
  \footnotesize
  \caption{Năm điều kiện, mỗi điều kiện là phần còn lại sau khi một chương đọc
  xong bài của nó. Cột phải là điều mà một phương pháp phải công bố thêm để thỏa
  điều kiện; điều kiện cuối là điều kiện duy nhất không đòi công bố thêm gì.}
  \label{tab:ch18-nam-dieu-kien}
  \begin{tabular}{@{}p{4cm}p{2cm}p{6cm}@{}}
    \toprule
    Điều kiện & Rút ra ở & Đòi công bố thêm gì \\
    \midrule
    Mô tả chỗ lệch, đừng xác nhận sự trùng khớp &
    mục~\ref{sec:ch13-tran-noi-gi} &
    Chỗ lệch nằm ở đâu trong không gian đầu vào, lệch theo hướng nào, và có rơi
    vào vùng người dùng quan tâm không \\
    \addlinespace
    Công bố tập can thiệp mà điểm số tương đối với nó &
    mục~\ref{sec:ch14-dieu-chua-giai-quyet} &
    Tập các phép thay và phép che mà điểm số được tính trên đó \\
    \addlinespace
    Báo phép ép cạnh phép đọc tên &
    mục~\ref{sec:ch15-doc-duoc-dieu-khien-duoc} &
    Kết quả của phép thử ép mỗi đơn vị, đặt cạnh cái tên đọc được từ đơn vị ấy \\
    \addlinespace
    Nói dụng cụ nằm dưới đã được kiểm định bằng gì &
    mục~\ref{sec:ch16-ba-cai-gia} &
    Với mỗi dụng cụ mà phép đo dựa lên, phép kiểm định của dụng cụ đó \\
    \addlinespace
    Nói nó hỏi câu nào trong hai câu &
    mục~\ref{sec:ch17-chieu-cua-phep-chuyen} &
    Không gì cả; chỉ đòi gọi đúng tên quan hệ đang đo \\
    \bottomrule
  \end{tabular}
\end{table}

\index{dụng cụ đo tương lai!điều kiện thứ nhất đến từ một định lý}%
Điều kiện thứ nhất khác bốn điều kiện kia ở chỗ nó đến từ một định lý chứ không
từ một phép đo. Định lý~\ref{thm:ch13-khoang-cach-do-phuc-tap} nói trước rằng một
lời giải thích đọc được không trùng khớp với một mô hình phức tạp, nên một dụng cụ
đo được giao việc xác nhận sự trùng khớp ấy là một dụng cụ được giao một việc đã
biết là không có. Việc còn lại cho nó là mô tả chỗ lệch, và cột phải của
bảng nói mô tả gồm ba phần nào.

Ba điều kiện tiếp theo có chung một hình dạng: cả ba đòi một phương pháp công bố
thêm một thứ mà hiện nay không ai bắt phải công bố. Điều kiện thứ hai đòi tập can
thiệp, và mục~\ref{sec:ch14-ground-truth-ke-tan-cong} cho thấy hậu quả của việc
không đòi, vì \tn{phép dựng giàn}{scaffolding} chạy được chính nhờ chỗ không ai
phải nói tập ấy là tập nào. Điều kiện thứ ba đòi phép ép đi kèm phép đọc tên, và con số đằng sau nó là tỉ lệ
những đơn vị mà cả hai phép cùng đạt: 18.84 phần trăm trên một mô hình, và
mục~\ref{sec:ch16-mot-cai-ten} đã đo lại trên ba mô hình để thấy tỉ lệ ấy đổi
theo mô hình và đổi theo ngưỡng, nên điều kiện là báo cả hai kết quả chứ không
phải giữ lấy một con số. Điều kiện thứ tư đòi khai ra dụng cụ nằm dưới, và
mục~\ref{sec:ch16-duong-thu-ba} là chỗ mà cả một lối thoát lẫn phép sửa của nó
cùng hụt điều kiện ấy.

\index{dụng cụ đo tương lai!điều kiện thứ năm không đòi công bố gì}%
Điều kiện thứ năm là điều kiện rẻ nhất và nó không đòi công bố thêm thứ gì. Nó
chỉ đòi một dụng cụ nói rõ nó hỏi câu nào trong hai câu mà
mục~\ref{sec:ch09-dinh-nghia-lai} tách ra: đường mà mô hình đã đi, tức độ trung
thực, hay ảnh hưởng nhân quả của lời giải thích lên đáp án, tức
\tn{mức quan trọng}{importance}.
Chương~\ref{ch:do-trung-thuc-chuoi-suy-luan} gặp ba bài in một cái tên cho cả hai
quan hệ, nên điều kiện ấy rẻ mà không tầm thường.

\index{dụng cụ đo tương lai!bài siêu đánh giá thỏa mấy điều kiện}%
Có một phép thử sẵn cho cả năm điều kiện, và mục~\ref{sec:ch18-ba-bai-le-ra-da-dong}
vừa đưa nó tới: đem năm điều kiện ra hỏi bài siêu đánh giá, là bài gần nhất với
một dụng cụ kiểm định chỉ số mà lĩnh vực có. Bài ấy thỏa điều kiện thứ tư một
cách mà không bài nào khác trong sách thỏa, vì đối tượng của nó \emph{là} dụng cụ
nằm dưới và nó nói thẳng rằng dụng cụ ấy chưa được kiểm định về tính hợp lệ. Điều
kiện thứ năm thì nó cũng thỏa, vì nó tuyên bố rõ nó đo độ tin cậy chứ không đo
tính hợp lệ. Ba điều kiện còn lại không áp cho nó, vì cả ba nói về một dụng cụ
chấm điểm lời giải thích, còn nó chấm điểm dụng cụ.

Nghĩa của phép thử ấy đáng ghi. Cái mà lĩnh vực đã dựng được thỏa đúng hai điều
kiện của sách, và cả hai đều là điều kiện về sự trung thực trong lời khai chứ
không phải về phép đo. Ba điều kiện đòi công bố thêm số liệu thì chưa dụng cụ nào
thỏa, vì chưa có dụng cụ nào ở tầng ấy để mà thỏa.

\index{dụng cụ đo tương lai!năm điều kiện|)}%
Năm điều kiện này không cộng lại thành một thiết kế, và mục này không giả vờ rằng
chúng làm được việc ấy. Chúng là năm ràng buộc mà năm chương khác nhau rút ra từ
năm hướng khác nhau, và giá trị của chúng nằm ở chỗ chúng loại bớt: một dụng cụ
không thỏa điều kiện thứ nhất là một dụng cụ được giao việc mà một định lý đã cấm,
và một dụng cụ không thỏa điều kiện thứ năm là một dụng cụ không nói được nó đo
cái gì. Mục sau đi từ ràng buộc sang việc: những phép đo mà sách gặp trên đường
đi, chưa ai chạy, mà không đòi dựng gì mới.
